
\documentclass[12pt]{amsart}
\usepackage{geometry} % see geometry.pdf on how to lay out the page. There's lots.
\geometry{a4paper} % or letter or a5paper or ... etc
\usepackage[T1]{fontenc}
\usepackage[latin9]{inputenc}
\usepackage{amsmath}
\usepackage{amsaddr}
\usepackage{hyperref}
\usepackage{dirtytalk}
\usepackage{float}
\usepackage{listings}
\usepackage{color}
\usepackage{tikz}
\usepackage{enumitem}

 
\definecolor{codegreen}{rgb}{0,0.6,0}
\definecolor{codegray}{rgb}{0.5,0.5,0.5}
\definecolor{stringcolor}{rgb}{0.7,0.23,0.36}
\definecolor{backcolour}{rgb}{0.95,0.95,0.92}
\definecolor{keycolor}{rgb}{0.007,0.01,1.0}
\definecolor{itemcolor}{rgb}{0.01,0.0,0.49}
 
\lstdefinestyle{mystyle}{
    %backgroundcolor=\color{backcolour},   
    commentstyle=\color{codegreen},
    keywordstyle=\color{keycolor},
    numberstyle=\tiny\color{codegray},
    stringstyle=\color{stringcolor},
    basicstyle=\footnotesize,
    breakatwhitespace=false,         
    breaklines=true,                 
    captionpos=b,                    
    keepspaces=true,                 
    numbers=left,                    
    numbersep=5pt,                  
    showspaces=false,                
    showstringspaces=false,
    showtabs=false,                  
    tabsize=2
}
 
\lstset{style=mystyle}

\lstdefinelanguage{Swift}{
  keywords={associatedtype, class, deinit, enum, extension, func, import, init, inout, internal, let, operator, private, protocol, public, static, struct, subscript, typealias, var, break, case, continue, default, defer, do, else, fallthrough, for, guard, if, in, repeat, return, switch, where, while, as, catch, dynamicType, false, is, nil, rethrows, super, self, Self, throw, throws, true, try, associativity, convenience, dynamic, didSet, final, get, infix, indirect, lazy, left, mutating, none, nonmutating, optional, override, postfix, precedence, prefix, Protocol, required, right, set, Type, unowned, weak, willSet},
  ndkeywords={class, export, boolean, throw, implements, import, this},
  sensitive=false,
  comment=[l]{//},
  morecomment=[s]{/*}{*/},
  morestring=[b]',
  morestring=[b]"
}

\lstset{emph={Int,count,abs,repeating,Array}, emphstyle=\color{itemcolor}}


\title{Week 07}

\date{\today}

\lstset{style=mystyle}

%%% BEGIN DOCUMENT
\begin{document}
\maketitle


Name: 
\\

Collaborators (if any): 

\section{Tasks}

\subsection{Asymptotic Growth}
Sort all the functions below in increasing order of asymptotic (big-$O$) growth. If some have the
same asymptotic growth, then be sure to indicate that. Note: lg means base 2. 
\begin{enumerate}
\item $5n$
\item $4$ lg $n$
\item $4$ lg lg $n$
\item $n^4$
\item $n^{1/2}$ $lg^4$ $n$
\item $(lg \: n)^{5\:lg\:n}$
\item $n^{lg\:n}$
\item $5^n$
\item $4^{n^4}$
\item $4^{4^n}$
\item $5^{5^n}$
\item $5^{5n}$
\item $n^{n^{1/5}}$
\item $n^{n/4}$
\item ${n/4}^{n/4}$
\end{enumerate}

\section{Exercises/Problems}

\subsection{Solve the Recurrences} 
Give asymptotic upper and lower bounds for $T(n)$ in each of the following recurrences. Assume
that $T(n)$ is constant for $n \leq 2$. Make your bounds as tight as possible, and justify your answers. 
If you feel unsure about this, brush up on The Master Method.
\begin{enumerate}[label=(\alph*)]
\item $T(n) = 4T(n/4) + 5n$
\item $ T(n) = 4T(n/5) + 5n$
\item $T(n) = 5T(n/4) + 4n$
\item $T(n) = 25T(n/5) + n^2$
\item $T(n) = 4T(n/5) + lg \: n$
\item $T(n) = 4T(n/5) + lg^5 \: n \sqrt n$
\item $T(n) = 4T(\sqrt n) + lg^5 \: n$
\item $T(n) = 4T(\sqrt n) + lg^2 \: n$
\item $T(n) = T(\sqrt n) + 5$
\item $ T(n) = T(n/2) + 2T(n/5) + T(n/10) + 4n$
\end{enumerate} 

\subsection{Franchise Location} 
Choco Donuts, a new chocolate milk and donuts restaurant chain, wants to maximize their total profit and build restaurants on many street corners. 
\\

The road network is an undirected graph $G = (V, E)$, where the potential restaurant
sites are the vertices of the graph. Each vertex $a$ has a nonnegative integer value $p_a$, which describes the potential profit of site $a$. Two restaurants cannot be built on adjacent vertices (to avoid self-competition). You are tasked with designing an algorithm that outputs the chosen set $U \subseteq V$ of sites that maximizes the total profit $\sum _{a \epsilon U} \: p_a$.
\\

For parts (a)-(c), suppose that the street network $G$ is acyclic, i.e., a tree.
\begin{enumerate}[label=(\alph*)]
\item Consider the following greedy restaurant-placement algorithm: Choose the highest-profit vertex $a_0$ in the tree (breaking ties according to some order on vertex names) and put it into U. Remove $a_0$ from further consideration, along with all of its neighbors in $G$. Repeat until no further vertices remain.
\\

Give a counterexample to show that this algorithm does not always give a restaurant
placement with the maximum profit. 
\\

\item Give an efficient algorithm to determine a placement with maximum profit.
\\

\item Suppose that, in the absence of good market research, Choco Donuts decides that all sites are
equally good, so the goal is simply to design a restaurant placement with the largest number of locations. Give a simple greedy algorithm for this case, and prove its correctness. 
\\

\item Now suppose that the graph is arbitrary and not necessarily acyclic. Give the fastest correct algorithm you can for solving the problem. 
\end{enumerate}

\end{document}
